%=============================================================
%  ejemplosABCDE_guia.tex
%  Guía de Ejemplos (Nivel Básico) — A, B, C, D, E
%  Algoritmos y Programación en Python — Primer Semestre
%  Prof. Dr. Aboud Barsekh-Onji
%  Universidad Anáhuac México  |  Facultad de Ingeniería
%=============================================================
\documentclass[12pt,a4paper]{article}

%--- Codificación y fuentes (xelatex) ---
\usepackage{fontspec}
\usepackage[spanish]{babel}
\setmainfont{Latin Modern Roman}
\setmonofont[HyphenChar=None]{DejaVu Sans Mono}

%--- Parche babel-spanish: evita usar hyphenchar en fuentes mono ---
\makeatletter
\def\es@chf{}
\makeatother

%--- Geometría ---
\usepackage[top=2.5cm, bottom=2.5cm, left=2.5cm, right=2.5cm, headheight=15pt]{geometry}

%--- Colores y gráficos ---
\usepackage[dvipsnames,svgnames,table]{xcolor}
\usepackage{graphicx}
\usepackage{tikz}

%--- Encabezados y pie de página ---
\usepackage{fancyhdr}
\pagestyle{fancy}
\fancyhf{}
\rhead{\small Algoritmos y Programación en Python — Nivel Básico}
\lhead{\small Prof. Dr. Aboud Barsekh-Onji}
\cfoot{\thepage}
\renewcommand{\headrulewidth}{0.4pt}

%--- Secciones con color ---
\usepackage{titlesec}
\definecolor{anaDark}{RGB}{0,51,102}
\definecolor{anaAccent}{RGB}{180,0,0}
\definecolor{solColor}{RGB}{0,100,60}
\titleformat{\section}{\large\bfseries\color{anaDark}}{}{0em}{\thesection\quad}[\color{anaAccent}\titlerule]
\titleformat{\subsection}{\normalsize\bfseries\color{anaDark}}{}{0em}{\thesubsection\quad}

%--- Cajas de texto ---
\usepackage{tcolorbox}
\tcbuselibrary{listings,skins,breakable}

%--- Colores código Python ---
\definecolor{codegreen}{rgb}{0,0.6,0}
\definecolor{codegray}{rgb}{0.5,0.5,0.5}
\definecolor{codepurple}{rgb}{0.45,0,0.75}
\definecolor{backcolour}{rgb}{0.97,0.97,1.0}
\definecolor{keyblue}{rgb}{0.0,0.25,0.75}

%--- Estilo de salida (texto plano) ---
\usepackage{listings}
\lstdefinestyle{OutputStyle}{
  language={},
  basicstyle=\ttfamily\footnotesize,
  breaklines=true,
  keepspaces=true,
  numbers=none,
  frame=none,
  backgroundcolor=\color{white},
  xleftmargin=4pt,
}

%--- Estilo listings Python ---
\lstdefinestyle{PythonStyle}{
  language=Python,
  basicstyle=\ttfamily\footnotesize,
  keywordstyle=\color{keyblue}\bfseries,
  commentstyle=\color{codegreen}\itshape,
  stringstyle=\color{codepurple},
  numberstyle=\tiny\color{codegray},
  breakatwhitespace=false,
  breaklines=true,
  captionpos=b,
  keepspaces=true,
  numbers=left,
  numbersep=7pt,
  showspaces=false,
  showstringspaces=false,
  showtabs=false,
  tabsize=4,
  frame=single,
  framerule=0.5pt,
  rulecolor=\color{anaDark!40},
  backgroundcolor=\color{backcolour},
  xleftmargin=10pt,
  xrightmargin=5pt,
}
\lstset{style=PythonStyle}

%--- Cajas especiales ---
\newtcolorbox{objetivobox}[1]{
  colback=anaDark!5, colframe=anaDark,
  fonttitle=\bfseries\small, title={#1},
  breakable, left=6pt, right=6pt, top=4pt, bottom=4pt
}
\newtcolorbox{conceptobox}[1]{
  colback=anaAccent!5, colframe=anaAccent,
  fonttitle=\bfseries\small, title={#1},
  breakable, left=6pt, right=6pt, top=4pt, bottom=4pt
}
\newtcolorbox{enunciadobox}[1]{
  colback=orange!6, colframe=orange!70!black,
  fonttitle=\bfseries\small, title={#1},
  breakable, left=6pt, right=6pt, top=4pt, bottom=4pt
}
\newtcolorbox{solucionbanner}{
  colback=solColor!10, colframe=solColor,
  left=6pt, right=6pt, top=4pt, bottom=4pt
}
\newtcolorbox{salidabox}[1]{
  colback=gray!8, colframe=gray!50,
  fonttitle=\bfseries\small\color{gray!70!black}, title={#1},
  breakable, left=6pt, right=6pt, top=4pt, bottom=4pt
}
\newtcolorbox{notabox}{
  colback=yellow!8, colframe=orange!60,
  left=6pt, right=6pt, top=4pt, bottom=4pt
}

%--- Hiperlinks ---
\usepackage{hyperref}
\hypersetup{colorlinks=true,linkcolor=anaDark,urlcolor=anaAccent}

%--- Matemáticas ---
\usepackage{amsmath}

%=============================================================
\begin{document}

%--- Portada ---
\begin{titlepage}
  \begin{center}
    \vspace*{1.5cm}
    {\Large\bfseries\color{anaDark} Universidad Anáhuac México}\\[0.4cm]
    {\large Facultad de Ingeniería}\\[2.5cm]
    \begin{tikzpicture}
      \draw[anaDark,line width=3pt] (0,0) -- (12,0);
      \draw[anaAccent,line width=1.5pt] (0,-0.25) -- (12,-0.25);
    \end{tikzpicture}\\[0.8cm]
    {\Huge\bfseries\color{anaDark} Guía de Ejemplos}\\[0.4cm]
    {\Large\color{anaDark!70} Algoritmos y Programación en Python}\\[0.3cm]
    {\large\color{anaAccent} Nivel Básico — Primer Semestre}\\[0.4cm]
    \begin{tikzpicture}
      \draw[anaAccent,line width=1.5pt] (0,0) -- (12,0);
      \draw[anaDark,line width=3pt] (0,-0.25) -- (12,-0.25);
    \end{tikzpicture}\\[2cm]
    {\large Cinco ejemplos (A–E) con formato enunciado / solución}\\[0.5cm]
    \begin{tabular}{cl}
      \textcolor{anaAccent}{\textbullet} & \textbf{A:} Ciclos \texttt{for} anidados — patrones y combinaciones \\[4pt]
      \textcolor{anaAccent}{\textbullet} & \textbf{B:} Módulo \texttt{random} — \texttt{randint} y \texttt{choice} \\[4pt]
      \textcolor{anaAccent}{\textbullet} & \textbf{C:} Ordenamiento — \texttt{.sort()} y \texttt{sorted()} \\[4pt]
      \textcolor{anaAccent}{\textbullet} & \textbf{D:} Ciclos \texttt{while} y \texttt{while True} \\[4pt]
      \textcolor{anaAccent}{\textbullet} & \textbf{E:} Sentencia \texttt{match}/\texttt{case} básica \\[4pt]
    \end{tabular}\\[2.5cm]
    {\large\bfseries Prof. DSc. Aboud Barsekh-Onji}\\[0.3cm]
    {\normalsize Facultad de Ingeniería | Universidad Anáhuac México}\\[1.5cm]
    {\normalsize Ciclo escolar 2025-2026}
  \end{center}
\end{titlepage}

%--- Índice ---
\tableofcontents
\newpage

%=============================================================
%  EJEMPLO A — ENUNCIADO
%=============================================================
\section{Ejemplo A — Ciclos \texttt{for} anidados}
\label{sec:ejemploA}

\begin{objetivobox}{Objetivo}
Practicar el uso de dos ciclos \texttt{for} anidados para imprimir
patrones en pantalla y generar combinaciones a partir de dos listas.
\end{objetivobox}

\vspace{0.6em}
\begin{conceptobox}{Conceptos necesarios}
\begin{itemize}
  \item \texttt{for variable in range(n):} — repite el bloque \texttt{n} veces.
  \item \textbf{Ciclo externo:} se mueve una vez por fila (como el horario del reloj).
  \item \textbf{Ciclo interno:} completa \emph{todas} sus iteraciones por cada
        iteración del externo (como el minutero del reloj).
  \item \texttt{print(..., end=' ')} — imprime sin saltar de línea.
  \item \texttt{print()} — imprime solo el salto de línea.
  \item \texttt{for elemento in lista:} — recorre cada elemento de una lista.
\end{itemize}
\end{conceptobox}

\vspace{0.8em}
\begin{enunciadobox}{Enunciado — \texttt{ejemploA.py}}
Escribe un programa con tres partes que usen ciclos \texttt{for} anidados:

\begin{enumerate}
  \item \textbf{Rectángulo de asteriscos.}
        Imprime en pantalla un rectángulo de \textbf{4 filas} por \textbf{5 columnas}
        de asteriscos.  Usa un ciclo \texttt{for} externo para las filas y
        uno interno para las columnas.  Al terminar cada fila, imprime
        un salto de línea con \texttt{print()}.

  \item \textbf{Tablas de multiplicar del 1, 2 y 3.}
        Usando dos ciclos \texttt{for} anidados, imprime las tablas del
        \textbf{1 al 3} (con multiplicadores del 1 al 5).
        El ciclo externo debe recorrer las tablas (1, 2, 3);
        el interno, los multiplicadores (1 al 5).
        Deja una línea en blanco entre cada tabla.

  \item \textbf{Combinaciones de frutas y colores.}
        Dadas las listas:
        \begin{center}
          \texttt{frutas = ['manzana', 'pera', 'uva']}\\
          \texttt{colores = ['rojo', 'verde']}
        \end{center}
        Usa dos ciclos \texttt{for} anidados para imprimir \emph{todas}
        las combinaciones posibles de fruta + color
        (p. ej. \texttt{manzana - rojo}, \texttt{manzana - verde}, etc.).
\end{enumerate}
\end{enunciadobox}

\newpage
%=============================================================
%  EJEMPLO A — SOLUCIÓN
%=============================================================
\begin{solucionbanner}
  {\large\bfseries\color{solColor} Solución — Ejemplo A: Ciclos \texttt{for} anidados}
  \hfill \texttt{ejemploA.py}
\end{solucionbanner}

\vspace{0.8em}
\lstinputlisting[caption={ejemploA.py — For anidados (básico)},
                 label={lst:ejemploA}]{ejemploA.py}

\vspace{0.8em}
\paragraph{Parte 1 — Rectángulo de asteriscos.}
El ciclo exterior recorre las filas; por cada fila, el ciclo interior
imprime 5 asteriscos con \texttt{end=' '} (sin saltar de línea).
Al finalizar el ciclo interior, el \texttt{print()} vacío produce
el salto de línea que termina la fila.  Total de iteraciones: $4 \times 5 = 20$.

\paragraph{Parte 2 — Tablas de multiplicar.}
El ciclo exterior toma el número de tabla (\texttt{tabla}: 1, 2, 3).
El ciclo interior hace el conteo de multiplicadores (\texttt{mult}: 1 a 5).
El \texttt{print()} entre tablas añade una línea en blanco para separar visualmente.

\paragraph{Parte 3 — Combinaciones.}
El ciclo exterior recorre \texttt{frutas} (3 elementos);
el interior recorre \texttt{colores} (2 elementos) completo para cada fruta.
Resultado: $3 \times 2 = 6$ combinaciones impresas.

\vspace{0.6em}
\begin{salidabox}{Salida esperada}
\begin{lstlisting}[style=OutputStyle]
* * * * *
* * * * *
* * * * *
* * * * *

  1 x 1 = 1
  1 x 2 = 2  ...  1 x 5 = 5

  2 x 1 = 2  ...  3 x 5 = 15

  manzana - rojo
  manzana - verde
  pera - rojo
  pera - verde
  uva - rojo
  uva - verde
\end{lstlisting}
\end{salidabox}

\newpage
%=============================================================
%  EJEMPLO B — ENUNCIADO
%=============================================================
\section{Ejemplo B — Módulo \texttt{random}}
\label{sec:ejemploB}

\begin{objetivobox}{Objetivo}
Aprender a generar números enteros aleatorios y a elegir elementos
al azar de una lista usando las dos funciones esenciales del módulo
\texttt{random}: \texttt{randint()} y \texttt{choice()}.
\end{objetivobox}

\vspace{0.6em}
\begin{conceptobox}{Conceptos necesarios}
\begin{itemize}
  \item \texttt{import random} — importar el módulo antes de usarlo.
  \item \texttt{random.randint(a, b)} — devuelve un entero aleatorio
        entre \texttt{a} y \texttt{b} inclusive (ambos extremos incluidos).
  \item \texttt{random.choice(lista)} — devuelve \emph{un} elemento
        elegido al azar de la lista; la lista no se modifica.
  \item Acumulador: variable que suma resultados dentro de un ciclo.
  \item Condicional \texttt{if/else} dentro de un \texttt{for}.
\end{itemize}
\end{conceptobox}

\vspace{0.5em}
\begin{notabox}
  \small\textbf{Nota:} los resultados varían en cada ejecución porque son aleatorios.
  Eso es normal y es precisamente el objetivo del módulo \texttt{random}.
\end{notabox}

\vspace{0.8em}
\begin{enunciadobox}{Enunciado — \texttt{ejemploB.py}}
Escribe un programa con tres partes usando \texttt{random}:

\begin{enumerate}
  \item \textbf{Dado de 6 caras.}
        Usa \texttt{random.randint(1, 6)} dentro de un ciclo \texttt{for}
        para simular \textbf{5 lanzamientos} de un dado.
        Imprime el número de lanzamiento y el resultado.

  \item \textbf{Sorteo de ganador.}
        Dada la lista:
        \begin{center}
          \texttt{participantes = ['Ana', 'Luis', 'Sofia', 'Carlos', 'Maria']}
        \end{center}
        Imprime la lista completa y luego usa \texttt{random.choice()}
        para elegir e imprimir \emph{un} ganador al azar.

  \item \textbf{Lanzamiento de moneda.}
        Usando \texttt{random.choice(['Cara', 'Sello'])} dentro de un
        ciclo \texttt{for} de 6 repeticiones:
        \begin{itemize}
          \item Imprime el resultado de cada lanzamiento.
          \item Cuenta cuántas veces salió 'Cara' y cuántas 'Sello'
                usando dos variables contadoras.
          \item Al final imprime el total de cada resultado.
        \end{itemize}
\end{enumerate}
\end{enunciadobox}

\newpage
%=============================================================
%  EJEMPLO B — SOLUCIÓN
%=============================================================
\begin{solucionbanner}
  {\large\bfseries\color{solColor} Solución — Ejemplo B: Módulo \texttt{random}}
  \hfill \texttt{ejemploB.py}
\end{solucionbanner}

\vspace{0.8em}
\lstinputlisting[caption={ejemploB.py — random.randint y random.choice},
                 label={lst:ejemploB}]{ejemploB.py}

\vspace{0.8em}
\paragraph{Parte 1 — Dado.}
\texttt{randint(1, 6)} genera un entero en $[1, 6]$ con igual probabilidad para
cada valor.  Se llama \emph{dentro} del ciclo para obtener un resultado distinto
en cada iteración.  Si se pusiera fuera del ciclo, todos los lanzamientos
mostrarían el mismo número.

\paragraph{Parte 2 — Sorteo.}
\texttt{random.choice(lista)} selecciona un elemento sin modificar la lista original.
No hace falta conocer los índices: la función se encarga de elegir uno al azar.

\paragraph{Parte 3 — Moneda y contadores.}
Las variables \texttt{caras} y \texttt{sellos} comienzan en 0 y se incrementan
con \texttt{+= 1} cada vez que el resultado coincide.
Al final, ambas suman 6 (el total de lanzamientos).

\vspace{0.6em}
\begin{salidabox}{Salida esperada (ejemplo — los valores cambian en cada ejecución)}
\begin{lstlisting}[style=OutputStyle]
  Lanzamiento 1: 4
  Lanzamiento 2: 1
  Lanzamiento 3: 6
  Lanzamiento 4: 3
  Lanzamiento 5: 5

  Participantes: ['Ana', 'Luis', 'Sofia', 'Carlos', 'Maria']
  El ganador es: Sofia

  Lanzamiento 1: Cara
  Lanzamiento 2: Sello
  ...
  Total caras:  3
  Total sellos: 3
\end{lstlisting}
\end{salidabox}

\newpage
%=============================================================
%  EJEMPLO C — ENUNCIADO
%=============================================================
\section{Ejemplo C — Ordenamiento de listas}
\label{sec:ejemploC}

\begin{objetivobox}{Objetivo}
Ordenar listas de números y palabras usando \texttt{.sort()} y
\texttt{sorted()}, y comprender la diferencia clave entre ambos métodos.
\end{objetivobox}

\vspace{0.6em}
\begin{conceptobox}{Conceptos necesarios}
\begin{itemize}
  \item \texttt{lista.sort()} — ordena la lista \textbf{en su lugar}
        (modifica el original).  No devuelve nada (\texttt{None}).
  \item \texttt{lista.sort(reverse=True)} — ordena de mayor a menor.
  \item \texttt{sorted(lista)} — devuelve una \textbf{nueva} lista ordenada;
        el original \textbf{no} cambia.
  \item Diferencia clave:
        \begin{center}
        \begin{tabular}{ll}
          \texttt{.sort()} & modifica la lista original \\
          \texttt{sorted()} & crea una lista nueva \\
        \end{tabular}
        \end{center}
\end{itemize}
\end{conceptobox}

\vspace{0.8em}
\begin{enunciadobox}{Enunciado — \texttt{ejemploC.py}}
Escribe un programa con tres partes sobre ordenamiento:

\begin{enumerate}
  \item \textbf{Ordenar calificaciones de forma ascendente con \texttt{.sort()}.}
        Dada la lista:
        \begin{center}
          \texttt{calificaciones = [75, 90, 60, 85, 70, 95, 65]}
        \end{center}
        Imprime la lista original, luego aplica \texttt{.sort()} e imprime
        la lista ordenada de menor a mayor.

  \item \textbf{Ordenar puntajes de forma descendente.}
        Dada la lista:
        \begin{center}
          \texttt{puntajes = [55, 88, 72, 99, 61, 80]}
        \end{center}
        Aplica \texttt{.sort(reverse=True)} para ordenarla de mayor a menor.
        Imprime el original y el resultado.

  \item \textbf{Ordenar nombres con \texttt{sorted()} sin modificar el original.}
        Dada la lista:
        \begin{center}
          \texttt{nombres = ['Carlos', 'Ana', 'Sofia', 'Luis', 'Maria', 'Beto']}
        \end{center}
        Usa \texttt{sorted(nombres)} para crear una nueva lista ordenada
        alfabéticamente.  Imprime: la lista original, la nueva lista ordenada,
        y confirma que el original \textbf{no cambió}.
\end{enumerate}
\end{enunciadobox}

\newpage
%=============================================================
%  EJEMPLO C — SOLUCIÓN
%=============================================================
\begin{solucionbanner}
  {\large\bfseries\color{solColor} Solución — Ejemplo C: Ordenamiento}
  \hfill \texttt{ejemploC.py}
\end{solucionbanner}

\vspace{0.8em}
\lstinputlisting[caption={ejemploC.py — .sort() y sorted()},
                 label={lst:ejemploC}]{ejemploC.py}

\vspace{0.8em}
\paragraph{Parte 1 — .sort() ascendente.}
\texttt{calificaciones.sort()} reorganiza los elementos dentro de la misma
lista.  Después de llamarlo, la variable \texttt{calificaciones} apunta
a la lista ya ordenada.  Imprimir antes y después de \texttt{.sort()} muestra
el cambio.

\paragraph{Parte 2 — .sort(reverse=True).}
El parámetro \texttt{reverse=True} invierte el criterio de ordenamiento
(de mayor a menor).  Por defecto \texttt{reverse=False}, es decir,
ascendente.

\paragraph{Parte 3 — sorted() sin modificar el original.}
\texttt{sorted(nombres)} devuelve una lista \emph{nueva} y deja
\texttt{nombres} intacto.  Por eso se necesita asignar el resultado:
\texttt{nombres\_ordenados = sorted(nombres)}.
La tercera impresión confirma que \texttt{nombres} no cambió.

\vspace{0.6em}
\begin{salidabox}{Salida esperada}
\begin{lstlisting}[style=OutputStyle]
  Original:    [75, 90, 60, 85, 70, 95, 65]
  Ascendente:  [60, 65, 70, 75, 85, 90, 95]

  Original:    [55, 88, 72, 99, 61, 80]
  Descendente: [99, 88, 80, 72, 61, 55]

  Original:    ['Carlos', 'Ana', 'Sofia', 'Luis', 'Maria', 'Beto']
  Ordenados:   ['Ana', 'Beto', 'Carlos', 'Luis', 'Maria', 'Sofia']
  Original ok: ['Carlos', 'Ana', 'Sofia', 'Luis', 'Maria', 'Beto']
\end{lstlisting}
\end{salidabox}

\newpage
%=============================================================
%  EJEMPLO D — ENUNCIADO
%=============================================================
\section{Ejemplo D — Ciclos \texttt{while} y \texttt{while True}}
\label{sec:ejemploD}

\begin{objetivobox}{Objetivo}
Entender y aplicar el ciclo \texttt{while} con una condición explícita
para contar y acumular, y el patrón \texttt{while True} con \texttt{break}
para controlar la salida del ciclo manualmente.
\end{objetivobox}

\vspace{0.6em}
\begin{conceptobox}{Conceptos necesarios}
\begin{itemize}
  \item \texttt{while condicion:} — repite el bloque \emph{mientras} la
        condición sea \texttt{True}.  La condición se evalúa al inicio
        de cada iteración.
  \item \textbf{Regla de oro:} el cuerpo del ciclo debe modificar la condición,
        o el ciclo se repetirá infinitamente.
  \item \texttt{variable += 1} — forma corta de \texttt{variable = variable + 1}.
  \item \textbf{Acumulador:} variable que suma valores dentro del ciclo
        (\texttt{suma += numero}).
  \item \texttt{while True:} — ciclo infinito; solo termina con \texttt{break}.
  \item \texttt{break} — sale inmediatamente del ciclo.
\end{itemize}
\end{conceptobox}

\vspace{0.8em}
\begin{enunciadobox}{Enunciado — \texttt{ejemploD.py}}
Escribe un programa con tres partes:

\begin{enumerate}
  \item \textbf{Contar del 1 al 5 con \texttt{while}.}
        Declara una variable \texttt{contador = 1}.
        Usa un ciclo \texttt{while} para imprimir el valor de \texttt{contador}
        mientras sea menor o igual a 5.
        No olvides incrementar el contador dentro del ciclo.
        Al salir del ciclo, imprime \texttt{'Ciclo terminado.'}.

  \item \textbf{Suma del 1 al 10 con \texttt{while}.}
        Declara \texttt{numero = 1} y \texttt{suma = 0}.
        Usa un ciclo \texttt{while} para sumar todos los enteros del 1 al 10.
        Al terminar, imprime el resultado total.
        (El resultado correcto es 55.)

  \item \textbf{\texttt{while True} con \texttt{break}.}
        Declara \texttt{intentos = 0}.
        Usa \texttt{while True} para repetir un bloque que:
        \begin{itemize}
          \item Incrementa \texttt{intentos} en 1.
          \item Imprime \texttt{'Intento numero X'}.
          \item Si \texttt{intentos} llega a 4, imprime un mensaje de límite
                y sale del ciclo con \texttt{break}.
        \end{itemize}
\end{enumerate}
\end{enunciadobox}

\newpage
%=============================================================
%  EJEMPLO D — SOLUCIÓN
%=============================================================
\begin{solucionbanner}
  {\large\bfseries\color{solColor} Solución — Ejemplo D: \texttt{while} y \texttt{while True}}
  \hfill \texttt{ejemploD.py}
\end{solucionbanner}

\vspace{0.8em}
\lstinputlisting[caption={ejemploD.py — while y while True},
                 label={lst:ejemploD}]{ejemploD.py}

\vspace{0.8em}
\paragraph{Parte 1 — Contador.}
La condición \texttt{while contador <= 5} se evalúa \emph{antes} de cada
iteración.  La línea \texttt{contador += 1} es imprescindible: sin ella
\texttt{contador} nunca cambiaría y el ciclo correría para siempre
(ciclo infinito).

\paragraph{Parte 2 — Acumulador.}
El patrón \texttt{suma += numero} es el acumulador: en cada iteración suma
el valor actual de \texttt{numero} a \texttt{suma}.  \texttt{numero} se
incrementa para avanzar hacia la condición de paro.  Cuando \texttt{numero}
llega a 11, la condición \texttt{numero <= 10} es \texttt{False} y el ciclo
termina.

\paragraph{Parte 3 — While True con break.}
\texttt{while True} crea un ciclo sin condición de paro natural.  El
\texttt{break} dentro del \texttt{if} es el único mecanismo de salida.
Este patrón es útil cuando la condición de paro se calcula \emph{dentro}
del cuerpo del ciclo.

\vspace{0.6em}
\begin{salidabox}{Salida esperada}
\begin{lstlisting}[style=OutputStyle]
  Contador: 1
  Contador: 2
  Contador: 3
  Contador: 4
  Contador: 5
  Ciclo terminado.

  Suma de 1 a 10 = 55

  Intento numero 1
  Intento numero 2
  Intento numero 3
  Intento numero 4
  Se alcanzo el limite. Saliendo...
\end{lstlisting}
\end{salidabox}

\newpage
%=============================================================
%  EJEMPLO E — ENUNCIADO
%=============================================================
\section{Ejemplo E — Sentencia \texttt{match}/\texttt{case}}
\label{sec:ejemploE}

\begin{objetivobox}{Objetivo}
Usar la sentencia \texttt{match}/\texttt{case} de Python 3.10+ para
evaluar una variable contra varios valores posibles, como alternativa
más legible a una cadena de \texttt{if/elif/else}.
\end{objetivobox}

\vspace{0.6em}
\begin{conceptobox}{Conceptos necesarios}
\begin{itemize}
  \item \texttt{match variable:} — evalúa \texttt{variable} y la compara
        con cada \texttt{case} en orden.
  \item \texttt{case valor:} — el bloque se ejecuta si \texttt{variable == valor}.
  \item \texttt{case \_:} — comodín: se ejecuta si \emph{ningún} caso
        anterior coincidió.  Equivale al \texttt{else} final.
  \item Solo el \textbf{primer} \texttt{case} que coincida se ejecuta;
        los demás se ignoran automáticamente (no es necesario \texttt{break}).
  \item Disponible desde \textbf{Python 3.10} en adelante.
\end{itemize}
\end{conceptobox}

\vspace{0.8em}
\begin{enunciadobox}{Enunciado — \texttt{ejemploE.py}}
Escribe un programa con tres partes que usen \texttt{match}/\texttt{case}:

\begin{enumerate}
  \item \textbf{Día de la semana.}
        Dado un número entero del 1 al 7, imprime el nombre del día
        correspondiente (\texttt{1} = Lunes, \ldots, \texttt{7} = Domingo).
        Para cualquier otro número, imprime \texttt{'Numero invalido'}.
        Prueba con los valores 1 al 8 (usando un ciclo \texttt{for}).

  \item \textbf{Semáforo.}
        Dada la lista \texttt{['rojo', 'amarillo', 'verde', 'azul']},
        recorre cada color con un ciclo \texttt{for} y usa
        \texttt{match}/\texttt{case} para imprimir la acción correspondiente:
        \begin{center}
        \begin{tabular}{ll}
          \texttt{rojo} & $\to$ \texttt{Detente} \\
          \texttt{amarillo} & $\to$ \texttt{Precaucion} \\
          \texttt{verde} & $\to$ \texttt{Avanza} \\
          cualquier otro & $\to$ \texttt{Color no reconocido} \\
        \end{tabular}
        \end{center}

  \item \textbf{Calificación en letra.}
        Dada la lista \texttt{['A', 'B', 'C', 'D', 'F', 'X']}, recorre
        cada letra con un ciclo \texttt{for} y usa \texttt{match}/\texttt{case}
        para imprimir la descripción de la calificación:
        A = Excelente (90--100), B = Bien (80--89), C = Suficiente (70--79),
        D = Mínimo (60--69), F = Reprobado ($<$60), cualquier otra = no reconocida.
\end{enumerate}
\end{enunciadobox}

\newpage
%=============================================================
%  EJEMPLO E — SOLUCIÓN
%=============================================================
\begin{solucionbanner}
  {\large\bfseries\color{solColor} Solución — Ejemplo E: \texttt{match}/\texttt{case}}
  \hfill \texttt{ejemploE.py}
\end{solucionbanner}

\vspace{0.8em}
\lstinputlisting[caption={ejemploE.py — match/case básico},
                 label={lst:ejemploE}]{ejemploE.py}

\vspace{0.8em}
\paragraph{Parte 1 — Días de la semana.}
Cada \texttt{case} compara \texttt{numero} con un entero literal.
El ciclo \texttt{for} prueba todos los valores del 1 al 8; el 8 no coincide
con ningún \texttt{case} concreto y cae al comodín \texttt{case \_:}.
No se necesita \texttt{break}: Python sale automáticamente al encontrar
el primer caso que coincide.

\paragraph{Parte 2 — Semáforo.}
La variable \texttt{accion} se asigna dentro del bloque \texttt{match} y
se imprime fuera.  El comodín \texttt{case \_:} captura colores no definidos
(aquí, \texttt{'azul'}).  Este patrón es más fácil de leer que cuatro
\texttt{if/elif} consecutivos.

\paragraph{Parte 3 — Calificación en letra.}
Los \texttt{case} comparan cadenas de texto de un carácter.
La \texttt{'X'} de prueba activa el comodín.  Notar que los \texttt{case}
distinguen mayúsculas y minúsculas: \texttt{'a'} y \texttt{'A'} son distintos.

\vspace{0.6em}
\begin{salidabox}{Salida esperada}
\begin{lstlisting}[style=OutputStyle]
  Dia 1: Lunes      Dia 5: Viernes
  Dia 2: Martes     Dia 6: Sabado
  Dia 3: Miercoles  Dia 7: Domingo
  Dia 4: Jueves     Dia 8: Numero invalido

  rojo       -> Detente
  amarillo   -> Precaucion
  verde      -> Avanza
  azul       -> Color no reconocido

  A: Excelente  (90 - 100)
  B: Bien       (80 - 89)
  F: Reprobado  (< 60)
  X: Letra no reconocida
\end{lstlisting}
\end{salidabox}

\newpage
%=============================================================
\section*{Resumen comparativo}
\addcontentsline{toc}{section}{Resumen comparativo}

\begin{table}[h!]
\centering
\renewcommand{\arraystretch}{1.5}
\begin{tabular}{|p{1.2cm}|p{3cm}|p{4cm}|p{5cm}|}
\hline
\rowcolor{anaDark!15}
\textbf{Ej.} & \textbf{Tema} & \textbf{Funciones/sintaxis clave} & \textbf{Lo que practica} \\
\hline
A & \texttt{for} anidados & \texttt{for}, \texttt{range()}, \texttt{print(end='')} & Patrones, tablas, combinaciones \\
\hline
B & \texttt{random} & \texttt{randint()}, \texttt{choice()} & Números y elementos aleatorios \\
\hline
C & Ordenamiento & \texttt{.sort()}, \texttt{.sort(reverse=True)}, \texttt{sorted()} & Ordenar números y palabras \\
\hline
D & \texttt{while} & \texttt{while cond:}, \texttt{while True:}, \texttt{break}, \texttt{+=} & Contador, acumulador, ciclo infinito \\
\hline
E & \texttt{match/case} & \texttt{match}, \texttt{case}, \texttt{case \_:} & Despacho por valor, comodín \\
\hline
\end{tabular}
\caption{Resumen de los cinco ejemplos de nivel básico}
\end{table}

\vspace{1.5em}
\begin{center}
  \small\color{gray}
  Universidad Anáhuac México — Facultad de Ingeniería\\
  Materia: Algoritmos y Programación en Python — Nivel Básico\\
  Prof. DSc. Aboud Barsekh-Onji — \texttt{aboud.barsekh@anahuac.mx}\\
  Ciclo 2025-2026
\end{center}

\end{document}
